\documentclass{report}
\usepackage[utf8]{inputenc}
\usepackage[english]{babel}
\usepackage{amsmath}
\usepackage{amssymb}

\begin{document}

\chapter*{Abstract}

The Boolean satisfiability problem (SAT) is studied in this work from
formal language theory. Within the line of research to which it belongs,
the satisfiability of a formula reduces to the emptiness of the
intersection of a regular language of interpretations with a grammar that
enforces consistency among the occurrences of each Boolean variable. With
less expressive grammars this reduction is limited to formulas with a
certain order of occurrences: a context-free grammar admits no crossings,
and with a simple range concatenation grammar (sRCG) the arity grows with
the formula. In an attempt to solve more formulas in polynomial time, the
problem is approached with range concatenation grammars (RCG): with their
non-linearity the consistency of more formulas is modeled while the arity
is kept minimal, as shown by an RCG that recognizes the consistency of
any SAT formula with arity one; but with that same non-linearity the
direct guarantee that the emptiness is polynomial is lost. A brute-force
algorithm that decides this emptiness is given: the strings of the
automaton are enumerated and recognized with the RCG; its cost is
exponential, so enumeration is not a feasible route, but it shows why the
problem is better studied on a single object. This object, the product
grammar, is constructed, and its size and the cost of constructing it are
characterized; on it the emptiness can be addressed: finding a polynomial
method for all formulas, which would imply $P = NP$; proving that none
exists; or finding one only by restricting the modeling to a subclass.
Although its size is not polynomial, it remains an object of study:
removing its non-productive clauses and analyzing the cost of deciding
the emptiness are proposed as future work.

\vspace{1em}
\noindent\textbf{Keywords:} boolean satisfiability, range concatenation
grammars, emptiness of the intersection, finite regular language,
non-linearity.

\end{document}
